\plannedchapter{进程、虚拟内存和地址空间}
{本章讲清楚用户态程序看到的地址空间如何由操作系统和硬件共同实现，解释为什么指针是虚拟地址，为什么 page、TLB 和 mmap 会影响性能。}
{\item 进程和线程的资源边界是什么？
\item 虚拟地址如何转换成物理地址？
\item ASLR、page fault、mmap、copy-on-write 对程序有什么影响？}
{\item 页表、TLB、page fault 的简化实现模型。
\item 用户态/内核态切换和系统调用的边界。
\item 地址空间布局与 text/data/bss/heap/stack/mmap regions 的关系。}
{打印不同对象地址；查看 \code{/proc/self/maps}；用大数组触发 page fault；比较首次访问和重复访问的时间差异。}
